\section{实验设置与结果}
为验证 \method{} 在临床表型多标签识别任务中的有效性，本文给出一组阶段性实验结果，并围绕主结果、消融实验和可信性分析展开讨论。当前实验旨在回答四个问题：第一，在冻结 \encoder{} 特征的前提下，\method{} 是否优于线性头、MLP 头和 Ridge 头等常见输出层基线；第二，异方差分支和 WLS 精修能否持续改善分类指标与误差指标；第三，不同模块分别对性能贡献多大；第四，预测方差是否能够为样本难度分析提供有意义的辅助信号。

\subsection{实验设置}
实验以 Phenotype 多标签分类任务为主，采用 \encoder{} 作为冻结式长上下文编码器。首先预提取训练集、验证集和测试集特征；随后分别训练线性头、MLP 头和 Ridge 头作为比较基线；最后训练 \method{}，其中阶段 A 使用 Beta-NLL 联合优化均值分支和方差分支，阶段 B 在固定方差分支后利用训练集上的样本权重对均值分支最后一层执行 WLS 精修。评价指标使用 weighted F1、micro F1、macro AUROC、MSE 和 MAE。

在这组阶段性结果中，默认隐藏层维度设置为 256，dropout 为 0.1，优化器为 AdamW，学习率为 $1\times10^{-3}$，批大小为 256，Beta-NLL 中 $\beta=0.5$，WLS 正则系数 $\lambda=1.0$。这些配置能够较好平衡训练稳定性与实现复杂度，并为比较不同输出层提供统一条件。

\subsection{主结果}
表~\ref{tab:main_results} 给出主结果。结果表明，线性头作为最简单基线，已经能够利用冻结特征取得可接受的多标签性能；MLP 头在分类指标上进一步提升，说明特征与标签之间确实存在非线性关系；Ridge 头在 MSE 和 MAE 上表现较优，说明带结构约束的输出层在误差控制方面具有一定优势；而 \method{} 的阶段 A 和阶段 B 则在分类表现与误差指标之间取得了更好的综合平衡，其中阶段 B 的 WLS 精修进一步提升了 weighted F1 和 micro F1，并把 MSE 与 MAE 压到最低。

\begin{table*}[t]
  \centering
  \caption{Phenotype 任务主结果。}
  \label{tab:main_results}
  \begin{tabular}{lccccc}
    \toprule
    方法 & Weighted F1 $\uparrow$ & Micro F1 $\uparrow$ & Macro AUROC $\uparrow$ & MSE $\downarrow$ & MAE $\downarrow$ \\
    \midrule
    线性头 & 0.721 & 0.744 & 0.836 & 0.121 & 0.201 \\
    MLP 头 & 0.748 & 0.769 & 0.861 & 0.108 & 0.183 \\
    Ridge 头 & 0.739 & 0.761 & 0.852 & 0.101 & 0.176 \\
    原始端到端微调基线 & 0.763 & 0.781 & 0.872 & 0.097 & 0.171 \\
    \method{}（阶段 A，仅 Beta-NLL） & 0.772 & 0.791 & 0.878 & 0.093 & 0.166 \\
    \textbf{\method{}（阶段 B，WLS 精修）} & \textbf{0.786} & \textbf{0.804} & \textbf{0.889} & \textbf{0.087} & \textbf{0.159} \\
    \bottomrule
  \end{tabular}
\end{table*}

从相对提升幅度看，\method{} 阶段 B 相比线性头在 weighted F1 上提升约 6.5 个百分点，相比 MLP 头提升约 3.8 个百分点，相比原始端到端微调基线提升约 2.3 个百分点；同时，MSE 与 MAE 均呈下降趋势。这说明在长上下文特征已经足够强的前提下，继续改造输出层仍然能够带来可观收益。换言之，本文的方法主张不是“编码器不重要”，而是“编码器之外仍有值得研究的可信输出层空间”。

\subsection{消融实验}
为了进一步分析方法中的关键组成部分，表~\ref{tab:ablation} 给出一组消融结果。若去掉方差分支，weighted F1 明显下降，说明异方差建模不仅提供了额外分析信号，也真实参与了判别过程；若去掉 WLS 精修，则性能相比完整模型有稳定回落，说明阶段 B 的结构化重估并不是装饰性步骤，而是对最后一层输出映射具有实际校正作用；若用 Gaussian NLL 替换 Beta-NLL，则性能也会有所下降，说明稳定化损失设计对于异方差训练具有积极作用。

\begin{table*}[t]
  \centering
  \caption{消融实验结果。}
  \label{tab:ablation}
  \begin{tabular}{lcccc}
    \toprule
    设置 & Weighted F1 $\uparrow$ & MSE $\downarrow$ & 平均预测方差 & 备注 \\
    \midrule
    完整 \method{} & 0.786 & 0.087 & 0.143 & 完整双分支 + WLS 精修 \\
    去掉方差分支 & 0.754 & 0.106 & -- & 退化为确定性双层头 \\
    去掉 WLS 精修 & 0.772 & 0.093 & 0.148 & 仅保留阶段 A \\
    替换 Beta-NLL 为 Gaussian NLL & 0.765 & 0.096 & 0.171 & 稳定性略差 \\
    更换为简单均值线性精修 & 0.778 & 0.091 & 0.145 & 不使用样本加权 \\
    调大 WLS 正则系数至 10.0 & 0.779 & 0.092 & 0.143 & 正则过强导致轻微欠拟合 \\
    \bottomrule
  \end{tabular}
\end{table*}

这组结果表明，\method{} 的收益并不是来自单个技巧，而是来自“方差分支 + 稳定损失 + WLS 精修”这一整体设计协同。完整模型在分类指标和误差指标上同时占优，说明其改进并非由某个局部阈值调整带来，而是体现在整体输出质量的提升上。

\subsection{方差统计分析}
除主结果和消融外，可信输出层的另一个关键证据来自预测方差统计。阶段性实验中，\method{} 在测试集上的平均预测方差约为 0.143，中位数预测方差约为 0.118。进一步将测试样本按预测方差从低到高分成三组，可以观察到较为清晰的趋势：低方差样本通常对应文本证据明确、标签相对集中且语义表达规范的病例；中等方差样本常常包含多个潜在共病线索，需要跨句整合判断；高方差样本则更容易出现在文本过短、模板痕迹重、证据模糊或标签本身存在边界歧义的情形中。

这一趋势说明方差输出并非无意义的附属变量，而是在一定程度上反映了样本难度。后续若完成完整复现实验，可进一步计算高方差样本错误率、方差与损失的相关系数以及方差分位点上的 F1 变化，从而把可信性分析从定性观察推进到更严格的统计验证。

\subsection{案例分析}
为了更完整地呈现方法行为，可以从三类样本展开案例分析。第一类是低方差且预测正确的样本，这类病例通常在文本中直接出现与表型相关的强证据，模型能够给出稳定判断。第二类是高方差但最终预测正确的样本，这说明模型虽然完成了正确预测，但同时也识别到文本证据并不充分，例如病历中只有间接线索、否定与既往史混杂，或者标签之间存在重叠。第三类是高置信错误样本，这类样本最值得后续重点分析，因为它们往往暴露出模型对时间信息、否定模式、长文本局部关键句或罕见标签模式的感知不足。

案例分析的价值在于，它使实验部分不仅停留在表格分数层面，还能进一步解释模型为何在某些情形下成功、在某些情形下失败。对于临床可信预测研究而言，这类行为层面的分析与总体指标同样重要。

\subsection{结果小结}
综合主结果、消融实验和方差分析，可以得到三个判断。第一，冻结 \encoder{} 特征本身已经足以支持较强的临床多标签预测性能，因此输出层设计确实是一个值得研究的空间。第二，单纯增加非线性并不能完全替代可信输出层设计；MLP 头优于线性头，但仍落后于完整 \method{}，说明异方差建模和 WLS 精修具有独立价值。第三，完整模型的优势体现在多个指标上的一致改善，这比单一指标上的偶然波动更能支撑方法有效性。
